% ============================ 04 小结 ============================

\section{小结、易错点与待核对项}

\begin{frame}{主要结论}
  \scriptsize
  \begin{block}{一、前向推理：网络预测的是“修正量”，不是轨迹}
    \scriptsize
    候选 $=$ 固定锚点 $+$ \textbf{网络预测的有界偏移与终端导数}；
    曲线由“当前状态 $+$ 终端状态”代入五次多项式边界值问题\textbf{解析求出}。
    在线只解被选中的一条；输入以\textbf{深度图为主}，状态次之，
    且状态会被旋转到 15 个基元系各一份。
  \end{block}
  \begin{block}{二、引导训练：监督来自“可求导的代价”，不是专家}
    \scriptsize
    评价器是手写的代价函数 $+$ 特权地图查询：
    平滑项与目标项给\textbf{闭式}梯度，避障项给\textbf{数值}梯度。
    自定义 autograd 层把 $\partial J/\partial\mathbf q$ 接进计算图，
    链式法则沿“解码链 $\to$ 网络权重”完成更新。
  \end{block}
  \begin{block}{三、两种输出、两条信号，一套权重}
    \scriptsize
    轨迹参数靠\textbf{外部代价梯度}改进（不经过预测评分）；
    评分靠\textbf{回归当前代价}学会排序；
    两者作用于共享骨干，是联合优化而非两个独立网络。
  \end{block}
\end{frame}

% ---------------------------------------------------------------
\begin{frame}{易混淆概念对照}
  \scriptsize
  \begin{center}
  \begin{tabular}{@{}p{0.33\textwidth}p{0.6\textwidth}@{}}
    \toprule
    \textbf{常见但不准确的理解} & \textbf{准确理解} \\
    \midrule
    15 个固定动作中选一个 & 固定的是锚点，终点是可预测的 \\
    15 个区域各运行搜索线程 & 共享网络并行产生多个空间位置的输出 \\
    只输出 2\,s 后的位置 & 还输出终端速度、加速度与评分 \\
    网络直接输出多项式系数 & 输出边界参数，系数由解析映射得到 \\
    在线生成全部曲线再查地图评分 & 参数与评分同时预测，只求解选中曲线 \\
    轨迹分支经“学出的评分”反传 & 主轨迹梯度来自外部代价评价器 \\
    无监督 $=$ 没有监督信号 & 无专家轨迹标签，但有代价与评分目标 \\
    终端限幅 $\Rightarrow$ 全程安全可执行 & 限幅与软代价改善表现，非全程硬约束证明 \\
    特权训练让网络知晓所有遮挡 & 部署仍受有限视野与部分可观测性限制 \\
    \bottomrule
  \end{tabular}
  \end{center}
  \vspace{0.2em}
  \begin{reading}
    \scriptsize
    由此可进一步推论：遇到严重遮挡、异常观测或训练分布之外的场景时，
    快速预测的评分可能失准。原文证明了特定场景中的有效性，
    \textbf{没有给出任意未知环境的完整安全保证}。
  \end{reading}
\end{frame}

% ---------------------------------------------------------------
\begin{frame}{待核对事项}
  \small
  \begin{columns}[t]
    \begin{column}{0.5\textwidth}
      \begin{block}{实现层面待确认}
        \footnotesize
        \begin{itemize}\tight
          \item 代价权重、归一化与时间采样间隔；
          \item 终端导数缩放的分量含义与坐标变换约定；
          \item 评分目标的符号、筛选阈值与比较方向；
          \item 自定义梯度、评分目标的停止梯度、聚合尺度；
          \item 重规划起点来自当前估计还是旧参考，
                以及拼接处的连续性。
        \end{itemize}
      \end{block}
    \end{column}
    \begin{column}{0.46\textwidth}
      \begin{block}{验证层面待确认}
        \footnotesize
        \begin{itemize}\tight
          \item 用简单障碍场景比较有限差分与解析/数值梯度；
          \item 轨迹表示换成其他参数化（如 MINCO）后，
                梯度的哪一段会失效；
          \item 论文与公开代码在评分符号、候选筛选、时域上的差异需要逐一记录。
        \end{itemize}
      \end{block}
      \begin{reading}
        \scriptsize
        作者仓库已有后续重构版本与不同轨迹表示的分支，
        参考时\textbf{应记录具体 commit}，
        避免把新版表示或部署指标写回 2024 年原文。
      \end{reading}
    \end{column}
  \end{columns}
\end{frame}

% ---------------------------------------------------------------
\begin{frame}{来源与说明}
  \small
  \begin{columns}[t]
    \begin{column}{0.5\textwidth}
      \begin{block}{本报告的依据}
        \footnotesize
        \begin{itemize}\tight
          \item 原文 PDF：\texttt{YOPO\_original.pdf}（8 页）；
                页码指该 PDF 页序，非期刊页码；
          \item 阅读笔记：\texttt{YOPO/papernote.md}（2026-10-08 版）；
          \item 原文配图：MinerU 提取结果中的图像文件；
          \item 笔记配图：由笔记作者绘制的教学示意图。
        \end{itemize}
      \end{block}
    \end{column}
    \begin{column}{0.46\textwidth}
      \begin{reading}[三类内容的标注]
        \footnotesize
        \begin{itemize}\tight
          \item \textbf{蓝色块}：原文可查证的结论，已注明图/表/式编号；
          \item \textbf{绿色块}：为便于理解构造的教学数字，
                不是论文实验数据；
          \item \textbf{橙色块}：本笔记的整理与推论，
                与原文陈述需区分。
        \end{itemize}
      \end{reading}
    \end{column}
  \end{columns}

  \vspace{0.4em}
  \begin{center}
    \Large\color{yopoblue}\bfseries 谢谢，欢迎讨论
  \end{center}
\end{frame}
